% part: intuitionistic-logic
% chapter: introduction
% section: bhk-interpretation

\documentclass[../../../include/open-logic-section]{subfiles}

\begin{document}

\olfileid[pt]{int}{int}{bhk}

\olsection{A interpretação de Brouwer--Heyting--Kolmogorov}

\begin{editorial}
  As provas de validade de proposições intuicionistas que usam a
  interpretação BHK são confusas; precisam ser melhor explicadas.
\end{editorial}

Há uma interpretação construtiva informal dos conectivos
intuicionistas, normalmente conhecida como interpretação de
Brouwer--Heyting--Kolmogorov. Ela usa a noção de ``construção'', que
pode ser entendida como uma prova construtiva. (Não usamos ``prova'' na
interpretação BHK para evitar confusão com a noção de !!a{derivation}
em um sistema formal de !!{derivation}.) Com base nessa noção intuitiva,
a interpretação BHK explica os significados dos conectivos
intuicionistas.

\begin{enumerate}
\item Supomos saber o que constitui uma construção de uma afirmação
  atômica.
\item Uma construção de $!A_1 \land !A_2$ é um par $\tuple{M_1, M_2}$,
  em que $M_1$ é uma construção de~$!A_1$ e $M_2$ é uma construção
  de~$!A_2$.
\item Uma construção de $!A_1 \lor !A_2$ é um par $\tuple{s, M}$,
  em que $s$ é~$1$ e $M$ é uma construção de~$!A_1$, ou $s$ é~$2$
  e $M$ é uma construção de~$!A_2$.
\item Uma construção de $!A \lif !B$ é uma função que converte uma
  construção de~$!A$ em uma construção de~$!B$.
\item Não há construção de~$\lfalse$ (o absurdo).
\item $\lnot !A$ é definida como sinônimo de $!A \lif \lfalse$. Isto é,
  uma construção de $\lnot !A$ é uma função que converte uma construção
  de~$!A$ em uma construção de~$\lfalse$.
\end{enumerate}

\begin{ex}
Considere, por exemplo, $\lnot \lfalse$. Uma construção dela é uma
função que, dada como entrada qualquer construção de $\lfalse$,
fornece como saída uma construção de~$\lfalse$. Evidentemente, a
função identidade~$\Id{}$ é uma construção desse tipo: dada uma
construção~$M$ de~$\lfalse$, $\Id{}(M) = M$ produz uma construção
de~$\lfalse$.
\end{ex}

Em termos gerais, $\lnot !A$ significa ``Uma construção de $!A$ é impossível''.

\begin{ex}
Provemos $!A \lif \lnot \lnot !A$ para qualquer proposição~$!A$, isto
é, $!A \lif ((!A \lif \lfalse) \lif \lfalse)$. A construção deve ser
uma função~$f$ que, dada uma construção~$M$ de $!A$, devolva uma
construção~$f(M)$ de $(!A \lif \lfalse) \lif \lfalse$. Eis como~$f$
constrói a construção de $(!A \lif \lfalse) \lif \lfalse$: precisamos
definir uma função $g$ que, ao receber como entrada uma construção~$h$
de $!A \lif \lfalse$, produza uma construção de~$\lfalse$. Podemos
definir $g$ assim: aplique a entrada~$h$ à construção~$M$ de $!A$ (que
recebemos anteriormente). Como a saída~$h(M)$ de $h$ é uma construção
de $\lfalse$, $f(M)(h) = h(M)$ é uma construção de~$\lfalse$ se $M$ é
uma construção de~$!A$.
\end{ex}

\begin{ex}
Apresentemos uma construção de $\lnot(!A \land \lnot !A)$, isto é,
$(!A \land (!A \lif \lfalse)) \lif \lfalse$. Trata-se de uma função~$f$
que, dada como entrada uma construção~$M$ de
$!A \land (!A \lif \lfalse)$, produz uma construção de~$\lfalse$. Uma
construção de uma conjunção $!B_1 \land !B_2$ é um par
$\tuple{N_1, N_2}$, em que $N_1$ é uma construção de~$!B_1$ e $N_2$ é
uma construção de~$!B_2$. Podemos definir funções $p_1$ e $p_2$ que,
a partir de uma construção de $!B_1 \land !B_2$, recuperam,
respectivamente, as construções de $!B_1$ e $!B_2$:
\begin{align*}
  p_1(\tuple{N_1, N_2}) &= N_1\\
  p_2(\tuple{N_1, N_2}) &= N_2
\end{align*}
Eis o que $f$ faz: primeiro, aplica $p_1$ à sua entrada~$M$. Isso
produz uma construção de~$!A$. Depois, aplica $p_2$ a~$M$, produzindo
uma construção de~$!A \lif \lfalse$. Essa construção, por sua vez, é
uma função~$p_2(M)$ que, se receber como entrada uma construção de~$!A$,
produz uma construção de~$\lfalse$. Em outras palavras, se aplicarmos
$p_2(M)$ a $p_1(M)$, obteremos uma construção de~$\lfalse$. Portanto,
podemos definir $f(M) = p_2(M)(p_1(M))$.
\end{ex}

\begin{ex}
Apresentemos uma construção de $((!A \land !B) \lif !C) \lif (!A \lif
(!B \lif !C))$, isto é, uma função $f$ que transforma uma construção~$g$
de $(!A \land !B) \lif !C$ em uma construção de $(!A \lif (!B \lif
!C))$. A construção $g$ é ela própria uma função (de construções de
$!A \land !B$ para construções de~$C$). A saída $f(g)$ é uma
função~$h_g$ que leva construções de $!A$ a funções que levam construções
de~$!B$ a construções de~$!C$.

Certo, isto é confuso. Precisamos construir certa função $h_g$, que
será a saída de $f$ para a entrada~$g$. A entrada de~$h_g$ é uma
construção~$M$ de~$!A$. A saída de $h_g(M)$ deve ser uma função~$k_M$
que leva construções~$N$ de~$!B$ a construções de~$!C$. Seja
$k_{g,M}(N) = g(\tuple{M, N})$. Recorde que $\tuple{M, N}$ é uma
construção de~$!A \land !B$. Assim, $k_{g,M}$ é uma construção de
$!B \lif !C$: ela leva construções~$N$ de~$!B$ a construções de~$!C$.
Agora, seja $h_g(M) = k_{g,M}$. Essa é uma função que leva construções~$M$
de~$!A$ a construções $k_{g,M}$ de $!B \lif !C$. Por fim, seja
$f(g) = h_g$. Essa é uma função que leva construções $g$ de
$(!A \land !B) \lif !C$ a construções de $!A \lif (!B \lif !C)$.
Ufa!{}
\end{ex}

A afirmação $!A \lor \lnot !A$ é chamada de lei do terceiro excluído.
Podemos prová-la para certos $!A$ específicos (por exemplo,
$\lfalse \lor \lnot \lfalse$), mas não em geral. Isso ocorre porque a
disjunção intuicionista exige uma construção de um dos disjuntos, mas
há afirmações que atualmente não podem ser provadas nem refutadas
(por exemplo, a conjectura de Goldbach). Contudo, tampouco se pode
refutar a lei do terceiro excluído: isto é, $\lnot \lnot (!A \lor
\lnot !A)$ vale.

\begin{ex}
Para provar $\lnot \lnot (!A \lor \lnot !A)$, precisamos de uma
função~$f$ que transforme uma construção de
$\lnot(!A \lor \lnot !A)$, isto é, de $(!A \lor (!A \lif \lfalse))
\lif \lfalse$, em uma construção de~$\lfalse$. Em outras palavras,
precisamos de uma função $f$ tal que $f(g)$ seja uma construção
de~$\lfalse$ se $g$ for uma construção de
$\lnot(!A \lor \lnot !A)$.

Suponha que $g$ seja uma construção de $\lnot(!A \lor \lnot !A)$,
isto é, uma função que transforma uma construção de
$!A \lor \lnot !A$ em uma construção de~$\lfalse$. Uma construção de
$!A \lor \lnot !A$ é um par $\tuple{s, M}$ no qual ou $s = 1$ e $M$
é uma construção de $!A$, ou $s = 2$ e $M$ é uma construção de
$\lnot !A$. Seja $h_1$ a função que leva uma construção~$M_1$ de $!A$
a uma construção de $!A \lor \lnot !A$: ela leva $M_1$ a
$\tuple{1, M_1}$. Seja $h_2$ a função que leva uma construção~$M_2$ de
$\lnot !A$ a uma construção de $!A \lor \lnot !A$: ela leva $M_2$ a
$\tuple{2, M_2}$.

Seja $k$ igual a $\comp{h_1}{g}$: trata-se de uma função que, dada uma
construção de~$!A$, devolve uma construção de~$\lfalse$, isto é, uma
construção de~$!A \lif \lfalse$ ou $\lnot !A$. Agora, seja $l$ igual a
$\comp{h_2}{g}$. Trata-se de uma função que, dada uma construção de
$\lnot !A$, fornece uma construção de~$\lfalse$. Como $k$ é uma
construção de $\lnot !A$, $l(k)$ é uma construção de~$\lfalse$.

Em conjunto, descrevemos como transformar uma construção~$g$ de
$\lnot(!A \lor \lnot !A)$ em uma construção de~$\lfalse$, isto é, a
função $f$ que leva uma construção~$g$ de
$\lnot(!A \lor \lnot !A)$ à construção $l(k)$ de~$\lfalse$ é uma
construção de $\lnot\lnot(!A \lor \lnot !A)$.
\end{ex}

Como se pode ver, usar a interpretação BHK para mostrar a validade
intuicionista de !!{formula}s rapidamente se torna trabalhoso e
confuso. Felizmente, há sistemas de !!{derivation} melhores para a
lógica intuicionista, bem como interpretações semânticas mais precisas.
\end{document}
